% Input fragment: requires amsmath, amssymb, amsthm and theorem/lemma/corollary environments.
% All statements and proofs below were independently derived for this packet.

\section{Positive Fourier annihilation with logarithmic support}
\label{sec:annihilation}

Our Fourier convention for a finite measure on the real line is
\[
 \widehat\mu(z)=\int_{\mathbb R}e^{izt}\,d\mu(t).
\]
If $\mu$ is supported on $[-H,H]$, this is entire. A zero of order
at least $r$ at $\lambda$ means
$\widehat\mu^{(j)}(\lambda)=0$ for $0\le j<r$.
For a finite symmetric multiset $\Lambda\subset\mathbb R\setminus\{0\}$,
write $m_\Lambda(\lambda)$ for its multiplicity. Symmetric means that
$m_\Lambda(-\lambda)=m_\Lambda(\lambda)$.

\subsection{A uniformly positive interpolating density}

\begin{lemma}[Separated frequencies]\label{lem:separated-filter}
Let $\Gamma\subset\mathbb R\setminus\{0\}$ be a finite symmetric set,
and suppose that the distinct points of $\Gamma\cup\{0\}$ are at least
$\delta>0$ apart. Put $A=4\pi/\delta$ and
\[
 g_A(t)=\frac1A\left(1-\frac{|t|}{A}\right)_+,
 \qquad
 f_A(z)=\left(\frac{\sin(Az/2)}{Az/2}\right)^2.
\]
There is an even probability density $w_\Gamma$ on $[-A,A]$ such that
\begin{equation}\label{eq:positive-envelope}
 \frac56 g_A(t)\le w_\Gamma(t)\le\frac65 g_A(t),
 \qquad
 \widehat w_\Gamma(\lambda)=0\quad(\lambda\in\Gamma).
\end{equation}
It is specified by a real linear system with a matrix whose induced
$\ell^1$ distance from the identity is at most $1/12$.
\end{lemma}

\begin{proof}
The density $g_A$ is the convolution of two uniform densities on
$[-A/2,A/2]$, so its Fourier transform is $f_A$. In particular, for
real $x\ne0$,
\[
 0\le f_A(x)\le\frac4{A^2x^2},\qquad f_A(0)=1.
\]
Index the coordinates by the distinct points of $\Gamma$, and define
\[
 M_{\lambda\kappa}=f_A(\lambda-\kappa),
 \qquad b_\lambda=f_A(\lambda).
\]
For each fixed $\kappa$, the $j$th point to its right is at distance
at least $j\delta$, and the same holds on its left. Consequently
\begin{align}
 \|M-I\|_1
 &\le \frac8{A^2\delta^2}\sum_{j\ge1}\frac1{j^2}
 =\frac{4\pi^2}{3A^2\delta^2}=\frac1{12},\label{eq:matrix-bound}\\
 \|b\|_1
 &\le \frac8{A^2\delta^2}\sum_{j\ge1}\frac1{j^2}
 =\frac1{12}.\label{eq:rhs-bound}
\end{align}
The second estimate uses separation from the additional point $0$.
The Neumann series makes $M$ invertible and gives
\[
 a=-M^{-1}b,\qquad
 \|a\|_1\le\frac{1/12}{1-1/12}=\frac1{11}.
\]
The data are real, so $a$ is real. Simultaneous replacement of
$\lambda,\kappa$ by $-\lambda,-\kappa$ preserves $M$ and $b$;
uniqueness gives $a_{-\lambda}=a_\lambda$. Therefore
\[
 Q(t)=1+\sum_{\lambda\in\Gamma}a_\lambda e^{-i\lambda t}
\]
is real and even, and $10/11\le Q(t)\le12/11$ for every $t$.
Its weighted integral
\[
 Z=\int_{\mathbb R}g_A(t)Q(t)\,dt
   =1+\sum_{\lambda\in\Gamma}a_\lambda f_A(\lambda)
\]
also lies between $10/11$ and $12/11$. Set
$w_\Gamma=g_AQ/Z$. It is an even probability density, and comparison
of these two bounds proves \eqref{eq:positive-envelope}. Finally,
\[
 \widehat w_\Gamma(\kappa)
 =\frac1Z\left(f_A(\kappa)+
       \sum_{\lambda\in\Gamma}a_\lambda
                      f_A(\kappa-\lambda)\right)=0
 \qquad(\kappa\in\Gamma),
\]
by $Ma=-b$.
\end{proof}

The constant $4\pi$ is convenient rather than optimized. The proof
uses finite matrices and a quantitative positivity margin; it invokes
no zero-separation conjecture for any $L$-function.

\subsection{Clustering, multiplicities, and uniform attenuation}

\begin{theorem}[A finite positive filter with an explicit support bound]
\label{thm:finite-filter}
Let $B>0$, and let $\Lambda$ be a finite symmetric multiset in
$[-B,B]\setminus\{0\}$. Choose $\delta>0$ such that every positive
element of $\Lambda$ is at least $\delta$. Define the local occupancy
\[
 K=\sup_{x\ge0}\sum_{\substack{\lambda>0\\x<\lambda\le x+\delta}}
                      m_\Lambda(\lambda).
\]
For every integer $N\ge0$, there is an even probability measure
$\mu_{\Lambda,N}$ supported on
\begin{equation}\label{eq:finite-support}
 [-H,H],\qquad H=\frac{4\pi K}{\delta}+\frac NB,
\end{equation}
whose transform satisfies
\begin{align}
 \operatorname{ord}_{z=\lambda}\widehat\mu_{\Lambda,N}(z)
   &\ge m_\Lambda(\lambda)
       &&(\lambda\in\Lambda),\label{eq:jet-zeros}\\
 |\widehat\mu_{\Lambda,N}(\omega)|
   &\le (\sin1)^N
       &&(\omega\in\mathbb R,\ |\omega|\ge B).\label{eq:tail-attenuation}
\end{align}
If $K+N>0$, the measure can be chosen absolutely continuous.
For $\Lambda$ empty, take $K=0$.
\end{theorem}

\begin{proof}
Process the positive occurrences of $\Lambda$ in nondecreasing order,
and color them greedily using $K$ colors. A color may be reused only
when its previous occurrence is at distance at least $\delta$.
At a current occurrence $\lambda$, every unavailable color has an
occurrence in $(\lambda-\delta,\lambda]$ that was processed earlier.
Together with the current occurrence there are at most $K$ such
occurrences. Thus at most $K-1$ colors are unavailable, and the
procedure succeeds. This argument includes repeated ordinates: equal
occurrences must receive different colors. The case
$\lambda-\delta<0$ causes no difficulty because no positive occurrence
is smaller than $\delta$.

For each nonempty color, reflect its positive set through the origin.
This gives a symmetric set $\Gamma_k$ with pairwise distance at least
$\delta$, also separated from $0$. In particular, the distance between
a positive and a negative point is at least $2\delta$. Apply
Lemma~\ref{lem:separated-filter} and convolve the resulting probability
measures. Their convolution $\nu$ is even, positive, and supported on
$[-4\pi K/\delta,4\pi K/\delta]$. Each occurrence of $\lambda$ belongs
to one color whose factor has a zero there. Since the Fourier transform
of a convolution is a product, its order at $\lambda$ is at least
$m_\Lambda(\lambda)$.

Let $u_B$ be the uniform probability measure on $[-1/B,1/B]$, and set
\[
 \mu_{\Lambda,N}=\nu*u_B^{*N},\qquad
 \widehat u_B(\omega)=\frac{\sin(\omega/B)}{\omega/B}.
\]
An empty convolution is the point mass at $0$. Convolution adds the
support radii and preserves the zero orders. On $1\le x\le\pi$,
$\sin x/x$ decreases from $\sin1$ to $0$, since the derivative's
numerator $x\cos x-\sin x$ has derivative $-x\sin x<0$. For
$x\ge\pi$, $|\sin x/x|\le1/\pi<\sin1$. Hence
\[
 \sup_{|x|\ge1}\left|\frac{\sin x}{x}\right|=\sin1<1.
\]
Every probability measure has Fourier transform of modulus at most
$1$ on the real axis. This proves \eqref{eq:tail-attenuation}, and the
absolute-continuity claim follows from the construction.
\end{proof}

There is a useful accuracy form of the same theorem. Set
$c_0=-\log(\sin1)>0$. Given $0<\varepsilon<1$, choose
$N=\lceil c_0^{-1}\log(1/\varepsilon)\rceil$. Then the tail is at most
$\varepsilon$ and
\[
 H\le\frac{4\pi K}{\delta}
       +\frac{\log(1/\varepsilon)}{c_0B}+\frac1B.
\]
Thus the support cost for exact low-frequency interpolation and the
additional cost for high-frequency damping are explicitly separated.

\begin{corollary}[Polynomial modes and frequency error]
\label{cor:jet-robustness}
Let $\mu$ be a compactly supported finite measure. If $\widehat\mu$ has a zero of order at least $r$ at a real
$\lambda$, then for every polynomial $P$ of degree less than $r$ and
every real $T$,
\[
 \int P(T+t)e^{i\lambda(T+t)}\,d\mu(t)=0.
\]
If moreover $\mu$ is a probability measure on $[-H,H]$, then for every
real $\eta$,
\begin{equation}\label{eq:frequency-robustness}
 |\widehat\mu(\lambda+\eta)|
 \le\min\left(1,\frac{(H|\eta|)^r}{r!}\right).
\end{equation}
\end{corollary}

\begin{proof}
Compact support permits differentiation under the integral:
$\widehat\mu^{(j)}(\lambda)=\int(it)^je^{i\lambda t}\,d\mu(t)$.
Expand $P(T+t)$ to obtain the first assertion. For the second, subtract
the first $r$ Taylor terms of $e^{i\eta t}$ inside the integral; their
integrals vanish. Taylor's integral remainder along the real axis is
bounded by $|\eta t|^r/r!$, because every derivative of $e^{ix}$ has
modulus one. Integrate and use $|t|\le H$. The additional bound $1$ is
the probability-measure estimate.
\end{proof}

\subsection{Application to Dirichlet zero ordinates}

Fix a finite family $\mathcal F$ of primitive Dirichlet characters
closed under complex conjugation, with the primitive trivial
character allowed. Let
$\Lambda_B$ be the multiset of all their nonzero zero ordinates
$\gamma=\Im\rho$ with $|\gamma|\le B$, counting multiplicity in each
$L$-function and adding multiplicities when zeros are shared.
Closure under conjugation makes $\Lambda_B$ symmetric. Zeros with
ordinate zero are excluded: no probability measure can have Fourier
transform zero at $0$. Under GRH these are precisely the central zeros.

The standard local zero bound \cite{Davenport} gives
\begin{equation}\label{eq:local-zero-bound}
 \sum_{\chi\in\mathcal F}
 \#\{\rho:L(\rho,\chi)=0,\ 0<\Re\rho<1,\ t<\Im\rho\le t+1\}
 \ll_{\mathcal F}\log(t+3),\qquad t\ge0,
\end{equation}
with multiplicities. Analyticity gives a positive minimum nonzero
absolute ordinate for this fixed finite family. Choose a fixed
$\delta\le1$ no larger than that minimum. Its occupancy $K_B$ in
Theorem~\ref{thm:finite-filter} is $O_{\mathcal F}(\log(B+3))$.
Taking $N=\lceil B\log B\rceil$, for $B\ge2$, proves the following.

\begin{theorem}[Logarithmic support for Dirichlet zero filters]
\label{thm:dirichlet-filter}
For the fixed conjugation-closed family $\mathcal F$, there are even probability densities $\mu_B$, $B\ge2$, with support radius
\[
 H_B=O_{\mathcal F}(\log(B+3)),
\]
which vanish at every prescribed nonzero zero ordinate up to $B$
to at least its aggregate multiplicity, and satisfy
\[
 \sup_{|\omega|\ge B}|\widehat\mu_B(\omega)|
 \le e^{-c_0B\log B},\qquad c_0=-\log(\sin1)>0.
\]
No assumption about simple zeros, distinct zeros of different
$L$-functions, or linear independence of ordinates is required.
\end{theorem}

This theorem about imaginary parts of zeros is unconditional: ordinates
are real numbers by definition. Dirichlet GRH is needed below to identify
the nontrivial-zero contributions to prime errors with pure Fourier modes.
The arithmetic input for the support estimate is the classical local zero
count, which follows from the standard zero-counting formula \cite{Davenport}. If the initial family is not closed
under conjugation, first adjoin its conjugates; the conclusion then
applies simultaneously to the original family.

\subsection{Why the logarithmic order cannot be reduced}

The following lower bound applies even to nonsymmetric probability
measures and even to a measure whose supporting interval is not
centered at the origin: translate the measure to center that interval,
which multiplies its Fourier transform by a nonvanishing exponential.

\begin{lemma}[Jensen support bound]\label{lem:jensen-support}
Let $\mu$ be a probability measure supported on $[-H,H]$. Suppose its
Fourier transform vanishes at every point of a finite symmetric
multiset $\Lambda\subset[-B,B]\setminus\{0\}$ to at least the
specified multiplicity. Then
\begin{equation}\label{eq:jensen-support}
 H\ge\frac{\pi}{2B}
       \sum_{\lambda\in\Lambda}\log\frac{B}{|\lambda|}
  =\frac\pi B\int_0^B\frac{N_\Lambda(t)}t\,dt,
\end{equation}
where $N_\Lambda(t)$ counts positive occurrences up to $t$.
\end{lemma}

\begin{proof}
Put $F(z)=\widehat\mu(z)$. Then $F(0)=1$ and
$|F(Re^{i\theta})|\le e^{HR|\sin\theta|}$. Jensen's formula, first
for radii without boundary zeros, gives
\[
 \sum_{F(z)=0,\ |z|<R}\log\frac R{|z|}
 =\frac1{2\pi}\int_0^{2\pi}\log|F(Re^{i\theta})|\,d\theta
 \le\frac{2HR}\pi.
\]
Keep just the prescribed zeros and let $R$ decrease to $B$ through
admissible radii. Zeros exactly at $B$ contribute zero in the limit.
This proves the first inequality. The equality is the elementary
identity
$\sum_{0<\lambda\le B}m_\Lambda(\lambda)\log(B/\lambda)
 =\int_0^BN_\Lambda(t)\,dt/t$, together with symmetry.
\end{proof}

For the zeta function, let $N(T)$ count positive ordinates with
multiplicity. The Riemann--von Mangoldt formula \cite{Montgomery,CGG} is
\[
 N(T)=\frac{T}{2\pi}\log\frac{T}{2\pi}
       -\frac{T}{2\pi}+O(\log(T+2)).
\]
Substitution into \eqref{eq:jensen-support} gives the explicit
asymptotic obstruction
\begin{equation}\label{eq:zeta-support-lower}
 H\ge\frac12\log B-\frac12\log(2\pi)-1
       +O\!\left(\frac{\log^2 B}{B}\right)
\end{equation}
for multiplicity-preserving cancellation of all zeta ordinates up to
$B$. Here the $O$ term means that the right-hand side can be decreased
by an absolute constant times $\log^2B/B$ to obtain an inequality.
More generally, for $d=|\mathcal F|$ with conductors $f_\chi$,
the corresponding aggregate-multiplicity lower bound is
\begin{equation}\label{eq:family-support-lower}
 H\ge\frac d2\log B+
       \frac12\sum_{\chi\in\mathcal F}\log\frac{f_\chi}{2\pi}
       -d+O_{\mathcal F}\!\left(\frac{\log^2 B}{B}\right).
\end{equation}
Indeed the sum of the positive zero-counting functions has the same
Riemann--von Mangoldt main terms, summed over the family. The
conjugation closure pairs any differences between positive and
negative ordinates. Thus the upper bound of
Theorem~\ref{thm:dirichlet-filter} has the optimal order $\log B$;
no claim is made that its leading constant is optimal.

\paragraph{Ordinary cancellation without prescribed multiplicity.}
For completeness, the logarithmic obstruction persists even when one
asks for only a simple Fourier zero at each distinct zeta ordinate.
Under RH, Montgomery's proved pair-sum formula \cite{Montgomery,CGG} implies
\[
 M_2(T):=\sum_{\substack{0<\gamma\le T\\\gamma\ {\text{distinct}}}}m_\gamma^2
 \le\left(\frac43+o(1)\right)N(T).
\]
This uses only the proved Fourier-support range of his theorem,
not the pair-correlation conjecture: the nonnegative squared-sinc
pair sum with parameter $0<a<1$ has asymptotic
$(a^{-1}+a/3)T\log T/(2\pi)$, and its equal-ordinate terms are exactly
$M_2(T)$; then let $a\uparrow1$ after $T\to\infty$.
For each integer $m\ge1$, $m-1\le(m^2-m)/2$. Hence, writing $D(T)$
for the number of distinct positive ordinates,
\[
 D(T)\ge\frac{3N(T)-M_2(T)}2
       \ge\left(\frac56-o(1)\right)N(T).
\]
Use $D(t)$ in Jensen's formula and integrate to obtain
\begin{equation}\label{eq:distinct-support-lower}
 H\ge\left(\frac5{12}-o(1)\right)\log B.
\end{equation}
Therefore even the weaker task of ordinary zero cancellation has
minimal support radius of order $\log B$. These constants are
convenient rigorous lower bounds, not asserted to be the best
constants available from the zero-density literature.

